\documentclass[12pt]{article}
\usepackage[T2A]{fontenc}
\usepackage[utf8]{inputenc}
\usepackage[russian]{babel}
\usepackage{amsmath}
\usepackage{booktabs}
\usepackage{geometry}
\geometry{margin=2cm}

\begin{document}

\begin{center}
\textbf{Лабораторная работа № 3}\\
\textbf{Создание электронной подписи в документе}\\
Вариант 11
\end{center}

\section{Цель работы}
Разработка процедур выработки и проверки электронной подписи (ЭП) сообщений на базе асимметричного криптографического алгоритма с применением функции хеширования.

\section{Постановка задачи}
\begin{enumerate}
    \item Выбрать алгоритм вычисления контрольной суммы (хэш-функции) по варианту.
    \item Разработать программную реализацию алгоритма создания и проверки ЭЦП.
    \item Подписать текстовое сообщение.
    \item Проверить правильность ЭП.
    \item Внести изменения в подпись или сообщение и убедиться, что подпись недействительна.
    \item Оформить результаты в виде отчёта.
\end{enumerate}

\section{Исходные данные}
По варианту 11 функция хэширования: \textbf{умножение кодов символов исходного текста по модулю \(2^{16}\)}. Алгоритм ЭЦП: \textbf{ГОСТ Р 34.10-94}. 

Выбранные параметры системы:
\begin{itemize}
    \item \(p = 23\) — простое число;
    \item \(q = 11\) — простое число, делитель \(p-1=22\);
    \item \(a = 2\), при этом \(a^q \bmod p = 2^{11} \bmod 23 = 1\);
    \item \(x = 5\) — секретный ключ;
    \item \(y = a^x \bmod p = 2^5 \bmod 23 = 9\) — открытый ключ;
    \item \(k = 7\) — сеансовый ключ;
    \item Сообщение \(M = \text{«SIGNATURE»}\).
\end{itemize}

\section{Алгоритм работы программы}
\begin{enumerate}
    \item Вычислить хэш сообщения \(h = H(M)\) по правилу варианта 11. Если \(h \bmod q = 0\), положить \(h = 1\).
    \item Сгенерировать случайное \(k \in [1, q-1]\).
    \item Вычислить \(r = a^k \bmod p\) и \(r_1 = r \bmod q\). Если \(r_1 = 0\), вернуться к п. 2.
    \item Вычислить \(s = (x \cdot r_1 + k \cdot h) \bmod q\). Если \(s = 0\), вернуться к п. 2.
    \item Подписью является пара \((r_1, s)\).
    \item Для проверки вычислить \(v = h^{q-2} \bmod q\).
    \item Вычислить \(z_1 = s \cdot v \bmod q\) и \(z_2 = (q - r_1) \cdot v \bmod q\).
    \item Вычислить \(u = (a^{z_1} \cdot y^{z_2} \bmod p) \bmod q\).
    \item Если \(u = r_1\), подпись верна, иначе — отвергается.
\end{enumerate}

\section{Результаты работы программы}
Хэш сообщения:
\[
h = \prod_{i} \text{ord}(M_i) \bmod 2^{16} = 51824.
\]
Проверка условия:
\[
h \bmod q = 51824 \bmod 11 = 3 \neq 0.
\]
Генерация подписи при \(k=7\):
\[
r = 2^7 \bmod 23 = 128 \bmod 23 = 13,
\]
\[
r_1 = 13 \bmod 11 = 2,
\]
\[
s = (5 \cdot 2 + 7 \cdot 51824) \bmod 11 = (10 + 362768) \bmod 11 = 9.
\]
Подпись: \((r_1, s) = (2, 9)\).

Проверка подписи:
\[
v = 3^{11-2} \bmod 11 = 3^9 \bmod 11 = 4,
\]
\[
z_1 = 9 \cdot 4 \bmod 11 = 3,
\]
\[
z_2 = (11 - 2) \cdot 4 \bmod 11 = 3,
\]
\[
u = (2^3 \cdot 9^3 \bmod 23) \bmod 11 = (8 \cdot 729 \bmod 23) \bmod 11 = 13 \bmod 11 = 2.
\]
Так как \(u = r_1\) (\(2 = 2\)), подпись \textbf{верна}.

При изменении сообщения на «SIGNATURF» хэш меняется, и проверка \(u = r_1\) не выполняется. Подпись отвергается.

\section{Анализ результатов}
Разработанный алгоритм корректно формирует и проверяет ЭЦП. При изменении исходного сообщения хэш-значение изменяется, что приводит к несовпадению \(u\) и \(r_1\). Подпись признается недействительной, что подтверждает защиту от подделки.

\section{Выводы}
В ходе работы изучен алгоритм ЭЦП ГОСТ Р 34.10-94. Реализована программная проверка подписи с использованием заданной хэш-функции. Подтверждено, что любое изменение данных приводит к недействительности подписи. Стойкость ЭП основана на сложности задачи дискретного логарифмирования.

\section{Ответы на контрольные вопросы}

\begin{enumerate}
    \item \textbf{Какие криптоалгоритмы используются для создания электронной подписи?} \\
    RSA, Эль-Гамаля, Шнорра, DSA, ГОСТ Р 34.10-94 и ГОСТ Р 34.10-2012. Они основаны на сложности факторизации или дискретного логарифмирования.

    \item \textbf{Что такое криптографическая хэш-функция, какими свойствами она должна обладать?} \\
    Это функция, отображающая данные произвольной длины в строку фиксированной длины. Свойства: однонаправленность (невозможно найти прообраз по хэшу), стойкость к коллизиям (сложно найти два разных сообщения с одинаковым хэшем) и лавинный эффект (малое изменение входа сильно меняет выход).

    \item \textbf{Как содержание сообщения влияет на электронную подпись?} \\
    Сообщение хэшируется, и подпись вычисляется на основе этого хэша. Любое изменение сообщения меняет хэш, что приводит к недействительности подписи.

    \item \textbf{Где используется ЭП?} \\
    В электронном документообороте, при подписании договоров, в банковских системах, для аутентификации программного обеспечения и защиты сетевых протоколов (например, HTTPS).

    \item \textbf{В каком случае электронная подпись при проверке отвергается?} \\
    Если хотя бы одно из условий \(0 < r_1 < q\), \(0 < s < q\) не выполнено, если хэш сообщения изменился, или если в результате вычислений \(u \neq r_1\).

    \item \textbf{От каких угроз информации защищает ЭП?} \\
    От подделки (изменения) документа, от отказа от авторства (невозможность заявить, что документ подписан не вами), а также обеспечивает аутентификацию источника.

    \item \textbf{Какие условия должны соблюдаться для параметра \(a\) в системе ЭП?} \\
    \(1 < a < p-1\), а также \(a^q \bmod p = 1\). Это гарантирует, что \(a\) является элементом подгруппы простого порядка \(q\).

    \item \textbf{Как открытый ключ \(y\) связан с секретным ключом \(x\)?} \\
    Соотношением \(y = a^x \bmod p\). Вычислить \(y\) по \(x\) легко, а найти \(x\) по \(y\) (задача дискретного логарифмирования) вычислительно сложно.

    \item \textbf{Почему контрольные суммы (например, CRC32) не подходят для использования в схемах ЭП?} \\
    Они не обладают криптостойкостью. CRC32 легко подделать, так как не является однонаправленной функцией и подвержена коллизиям, что позволяет злоумышленнику изменить документ, сохранив ту же контрольную сумму.

    \item \textbf{Как функция хеширования связана с защитой от изменений документа?} \\
    Хэш выступает как «цифровой отпечаток» документа. Подписывается именно хэш. Любое изменение документа меняет хэш, делая подпись невалидной.

    \item \textbf{Почему при внесении изменений в подписанное сообщение подпись становится недействительной?} \\
    Потому что при проверке заново вычисляется хэш от измененного сообщения. Новый хэш не совпадает с тем, который использовался при создании подписи, поэтому проверочное уравнение \(u = r_1\) не выполняется.

    \item \textbf{Почему в алгоритме ГОСТ Р34.10-94 длина подписи составляет 512 бит?} \\
    Подпись состоит из двух чисел \(r_1\) и \(s\), каждое из которых имеет размер до 256 бит (так как \(q\) имеет длину до 256 бит). Суммарно: \(256 + 256 = 512\) бит.

    \item \textbf{Какие шаги необходимо выполнить, чтобы подделать электронную подпись, и почему это вычислительно сложно?} \\
    Для подделки нужно вычислить секретный ключ \(x\) из уравнения \(y = a^x \bmod p\) или подобрать сеансовый ключ \(k\). Это требует решения задачи дискретного логарифмирования, для которой не существует эффективных алгоритмов.

    \item \textbf{Почему размеры параметров \(p\) и \(q\) строго регламентированы в ГОСТ?} \\
    Чтобы обеспечить требуемый уровень криптостойкости (не менее 256 бит для \(q\) и 1024 бит для \(p\) в старых версиях). Малые размеры позволят злоумышленнику факторизовать модуль или решить задачу дискретного логарифма за приемлемое время.

    \item \textbf{Почему в алгоритме требуется, чтобы параметр \(q\) был простым делителем \(p-1\)?} \\
    По теореме Лагранжа порядок подгруппы должен делить порядок группы. Если \(q\) — простой делитель \(p-1\), то существует подгруппа простого порядка \(q\). Это гарантирует, что вычисления в подгруппе не подвержены атакам на малые подгруппы, и обеспечивает высокую криптостойкость.
\end{enumerate}

\section*{Примечание}
Текст программы в отчёт не включён. Python-код находится в отдельном файле \texttt{eds.py}.

\end{document}